\documentclass[12pt,a4paper]{article}

%% -- Encoding & Language -------------------------------------
\usepackage[T1]{fontenc}
\usepackage[utf8]{inputenc}
\usepackage[english]{babel}

%% -- Mathematics ----------------------------------------------
\usepackage{mathtools,amssymb,amsthm}
\allowdisplaybreaks
\setlength{\emergencystretch}{3em}

%% -- Page Geometry --------------------------------------------
\usepackage[left=2.5cm,right=2cm,top=2.5cm,bottom=2.5cm]{geometry}

%% -- Tables & Lists -------------------------------------------
\usepackage{booktabs}
\usepackage{array}
\usepackage{enumitem}
\usepackage{listings}

%% -- Colors & Boxes -------------------------------------------
\usepackage{xcolor}
\usepackage[most]{tcolorbox}
\tcbuselibrary{skins,breakable}

%% -- Header/Footer --------------------------------------------
\usepackage{fancyhdr}
\usepackage{hyperref}
\hypersetup{colorlinks=true, linkcolor=blue!60!black, urlcolor=blue}

%% -- Color Definitions ----------------------------------------
\definecolor{kseblue}{RGB}{0,70,140}
\definecolor{kseyellow}{RGB}{255,185,0}
\definecolor{kselightblue}{RGB}{220,233,255}
\definecolor{kselightgreen}{RGB}{215,245,225}
\definecolor{kselightyellow}{RGB}{255,250,210}
\definecolor{kselightorange}{RGB}{255,238,210}
\definecolor{proofbg}{RGB}{240,240,255}
\definecolor{econgreen}{RGB}{0,110,60}
\definecolor{algobg}{RGB}{255,252,230}

%% -- Box Styles -----------------------------------------------
\newtcolorbox{defbox}[1]{
  enhanced, breakable,
  colback=kselightblue, colframe=kseblue,
  fonttitle=\bfseries\color{white},
  title={#1},
  attach boxed title to top left={yshift=-2mm,xshift=5mm},
  boxed title style={colback=kseblue, colframe=kseblue},
  left=5pt, right=5pt, top=6pt, bottom=5pt
}

\newtcolorbox{problembox}[1]{
  enhanced, breakable,
  colback=kselightyellow, colframe=kseyellow!80!black,
  fonttitle=\bfseries,
  title={#1},
  attach boxed title to top left={yshift=-2mm,xshift=5mm},
  boxed title style={colback=kseyellow!80!black,
    colframe=kseyellow!80!black, fontupper=\color{white}},
  left=5pt, right=5pt, top=6pt, bottom=5pt
}

\newtcolorbox{solutionbox}[1]{
  enhanced, breakable,
  colback=kselightgreen, colframe=econgreen,
  fonttitle=\bfseries\color{white},
  title={#1},
  attach boxed title to top left={yshift=-2mm,xshift=5mm},
  boxed title style={colback=econgreen},
  left=5pt, right=5pt, top=6pt, bottom=5pt
}

\newtcolorbox{algobox}[1]{
  enhanced, breakable,
  colback=algobg, colframe=orange!70!black,
  fonttitle=\bfseries,
  title={#1},
  attach boxed title to top left={yshift=-2mm,xshift=5mm},
  boxed title style={colback=orange!70!black, fontupper=\color{white}},
  left=5pt, right=5pt, top=6pt, bottom=5pt
}

\newtcolorbox{proofbox}[1]{
  enhanced, breakable,
  colback=proofbg, colframe=kseblue!60,
  fonttitle=\bfseries\itshape,
  title={#1},
  attach boxed title to top left={yshift=-2mm,xshift=5mm},
  boxed title style={colback=kseblue!60, fontupper=\color{white}},
  left=5pt, right=5pt, top=6pt, bottom=5pt
}

\newtcolorbox{notebox}{
  enhanced,
  colback=kselightorange, colframe=orange!60!black,
  left=5pt, right=5pt, top=4pt, bottom=4pt,
  fontupper=\small
}

\lstdefinestyle{rcode}{
  language=R,
  basicstyle=\ttfamily\small,
  breaklines=true,
  showstringspaces=false,
  frame=single,
  framerule=0.3pt
}

%% -- Header/Footer --------------------------------------------
\pagestyle{fancy}
\fancyhf{}
\rhead{\textcolor{kseblue}{\small Probability \& Statistics --- KSE}}
\lhead{\textcolor{kseblue}{\small Practice Worksheet (Lectures 3--4)}}
\cfoot{\thepage}
\renewcommand{\headrulewidth}{0.5pt}

%% -- Math Macros ----------------------------------------------
\newcommand{\E}{\mathbb{E}}
\newcommand{\Var}{\mathrm{Var}}
\newcommand{\Cov}{\mathrm{Cov}}
\newcommand{\N}{\mathcal{N}}
\renewcommand{\P}{\mathrm{P}}

\setlength{\parindent}{0pt}
\setlength{\parskip}{4pt}

\begin{document}

\begin{center}
  {\LARGE\bfseries\color{kseblue}
    Practice Worksheet in Probability \& Statistics}\\[6pt]
  {\large\bfseries for Economics Students}\\[3pt]
  {\normalsize\color{gray}
    Based on Lectures 3--4 \quad|\quad
    Rozora~I.V. \quad|\quad KSE, Kyiv, 2026}\\[4pt]
  {\small\color{gray}
    Topics:\;
    Normal distribution $\cdot$
    Joint distribution $\cdot$
    Independence $\cdot$
    Covariance $\cdot$
    Correlation $\cdot$
    Markov \& Chebyshev inequalities $\cdot$
    $3\sigma/6\sigma$ rules $\cdot$ LLN $\cdot$ CLT
  }
  \vspace{4pt}\hrule height 1.2pt\vspace{4pt}
\end{center}

\section*{\color{kseblue}Key Formulas \& Theory (Reference Card)}

\begin{defbox}{Normal Distribution\; $\xi \sim \N(\mu,\sigma^2)$}
\textbf{PDF:}\quad
\[
f(x)=\frac{1}{\sigma\sqrt{2\pi}}
\exp\!\Bigl(-\frac{(x-\mu)^2}{2\sigma^2}\Bigr),
\qquad
\E\xi=\mu,\quad \Var\xi=\sigma^2.
\]

\textbf{Z-standardisation:}\quad
$Z=\dfrac{\xi-\mu}{\sigma}\sim\N(0,1)$;\qquad
symmetry property:\; $\Phi(-a)=1-\Phi(a)$.

\textbf{Linear transform:}\quad
If $\xi\sim\N(\mu,\sigma^2)$ then $a\xi+b\sim\N(a\mu+b,\,a^2\sigma^2)$.

\textbf{Additivity (independent normals):}\quad
If $\xi_i\sim\N(\mu_i,\sigma_i^2)$ independent, then
\[
\sum_{i=1}^n\xi_i
\sim
\N\!\Bigl(\sum_{i=1}^n\mu_i,\;\sum_{i=1}^n\sigma_i^2\Bigr).
\]
\end{defbox}

\vspace{4pt}

\begin{defbox}{Joint Distribution \& Statistical Independence}
\textbf{Discrete joint PMF:}\;
$p(x,y)=\P(X=x,Y=y)$,\; $\sum_x\sum_y p(x,y)=1$.

\textbf{Marginals:}\;
$p_X(x)=\sum_y p(x,y)$ (sum over row);\quad
$p_Y(y)=\sum_x p(x,y)$ (sum over column).

\textbf{Continuous joint PDF:}\;
\[
\P((X,Y)\in A)=\iint_A f(x,y)\,dx\,dy,\qquad
f_X(x)=\int_{-\infty}^{\infty}f(x,y)\,dy.
\]

\textbf{Independence:}\;
\[
X\perp Y
\Longleftrightarrow
p(x,y)=p_X(x)\cdot p_Y(y)
\text{ for all }(x,y).
\]
\end{defbox}

\vspace{4pt}

\begin{defbox}{Covariance \& Pearson Correlation Coefficient}
\textbf{Covariance (computational formula):}\quad
$\Cov(X,Y)=\E[XY]-\E[X]\E[Y]$.

\textbf{Key properties:}
\begin{itemize}[noitemsep, topsep=2pt]
  \item $\Cov(X,X)=\Var X$
  \item $\Cov(aX+b,\,cY+d)=ac\,\Cov(X,Y)$
  \item $\Var(X+Y)=\Var X+\Var Y+2\Cov(X,Y)$
  \item $X\perp Y \Rightarrow \Cov(X,Y)=0$
        \textbf{(the converse is not generally true)}
\end{itemize}

\textbf{Pearson correlation:}\quad
\[
\rho_{X,Y}=\dfrac{\Cov(X,Y)}{\sigma_X\,\sigma_Y},\qquad
-1\le\rho\le1.
\]
\end{defbox}

\vspace{4pt}

\begin{defbox}{Markov's \& Chebyshev's Inequalities, $3\sigma$/$6\sigma$ Rules}
\textbf{Markov's Inequality}\; (requires $X\ge0$, $a>0$):
\[
  \P\{X>a\}\le\frac{\E X}{a}.
\]

\textbf{Chebyshev's Inequality}\; ($k>0$):
\[
  \P\{|X-\E X|>k\}\le\frac{\Var X}{k^2}.
\]
\end{defbox}

\vspace{4pt}
\begin{proofbox}{Remark: Proof of Chebyshev's Inequality (derived from Markov)}
\textbf{Theorem (Chebyshev).}\;
Let $X$ be a random variable with mean $\mu=\E X$ and finite variance
$\sigma^2=\Var X$. For any $k>0$:
\[
  \P\{|X-\mu|>k\}\le\frac{\sigma^2}{k^2}.
\]

\textbf{Proof.}
\begin{enumerate}[leftmargin=*,itemsep=2pt]
  \item Define $Y=(X-\mu)^2\ge0$, so $\E[Y]=\sigma^2$.
  \item $\{|X-\mu|>k\}=\{Y>k^2\}$.
  \item By Markov:
  \[
    \P\{Y>k^2\}\le\frac{\E[Y]}{k^2}=\frac{\sigma^2}{k^2}.
  \]
\end{enumerate}
\end{proofbox}

\vspace{4pt}

\begin{defbox}{Law of Large Numbers (LLN) \& Central Limit Theorem (CLT)}
\textbf{LLN (Chebyshev form):}
\[
\bar X_n=\frac{1}{n}\sum_{i=1}^n X_i\xrightarrow{p}\mu
\quad\text{as }n\to\infty.
\]

\textbf{CLT:}
\[
\frac{S_n-n\mu}{\sigma\sqrt{n}}\xrightarrow{d}\N(0,1),
\qquad S_n=X_1+\cdots+X_n.
\]
\end{defbox}

\bigskip
\section*{\color{kseblue}Problems from the Lecture}

\begin{problembox}{Problem 1}
An organization accepts members in the \textbf{top 2\%} of IQ scores.
Assume $X\sim\N(\mu=100,\;\sigma^2=16^2)$.
\end{problembox}

\begin{solutionbox}{Solution to Problem~1}
\[
x_{\min}=100+2.054\times16\approx132.9.
\]

\textbf{R code:}
\begin{lstlisting}[style=rcode]
mu <- 100
sd <- 16

# (a) Membership cutoff: top 2%
x_min <- qnorm(0.98, mean = mu, sd = sd)
x_min

# (b) Three random people from population all admitted
p_admit <- 0.02
p_all_three <- p_admit^3
p_all_three

# (c) Three random members all have IQ > 140
p_gt_140 <- 1 - pnorm(140, mean = mu, sd = sd)
p_cond <- p_gt_140 / p_admit
p_three_members <- p_cond^3
c(p_gt_140 = p_gt_140, p_cond = p_cond, p_three_members = p_three_members)
\end{lstlisting}
\end{solutionbox}

\vspace{6pt}

\begin{problembox}{Problem 2}
The lifespan of LED bulbs is normally distributed:
$X\sim\N(\mu=10{,}000,\;\sigma=500)$ hours.
\end{problembox}

\begin{solutionbox}{Solution to Problem~2}
\[
\begin{aligned}
\P(9000\le X\le11000)
&=\P\!\left(\frac{9000-10000}{500}\le Z\le\frac{11000-10000}{500}\right)\\
&=\P(-2\le Z\le2)=2\,\Phi_0(2)=2\times0.4772=0.9545.
\end{aligned}
\]
\textbf{R code:}
\begin{lstlisting}[style=rcode]
mu <- 10000
sd <- 500

p <- pnorm(11000, mean = mu, sd = sd) - pnorm(9000, mean = mu, sd = sd)
p
\end{lstlisting}
\end{solutionbox}

\vspace{6pt}

\begin{problembox}{Problem 3}
Two assets satisfy:
$\E[X]=5$, $\E[Y]=10$, $\E[XY]=52$, $\sigma_X=2$, $\sigma_Y=4$.
\end{problembox}

\begin{solutionbox}{Solution to Problem~3}
\[
\Cov(X,Y)=\E[XY]-\E[X]\E[Y]=52-5\times10=2,\quad
\rho=\frac{2}{2\times4}=0.25.
\]
\textbf{R code:}
\begin{lstlisting}[style=rcode]
EX <- 5
EY <- 10
EXY <- 52
sdX <- 2
sdY <- 4

covXY <- EXY - EX * EY
rho <- covXY / (sdX * sdY)
c(covXY = covXY, rho = rho)
\end{lstlisting}
\end{solutionbox}

\vspace{6pt}

\begin{problembox}{Problem 4}
A software firm tracks daily server crashes ($X$) and emergency hotfixes ($Y$).
\end{problembox}

\begin{solutionbox}{Solution to Problem~4}
\[
\begin{aligned}
\E[XY]
&=\sum_{x,y}xy\,p(x,y)\\
&=(1)(1)(0.25)+(1)(2)(0.10)+(2)(1)(0.05)+(2)(2)(0.05)\\
&=0.75.
\end{aligned}
\]
\textbf{R code:}
\begin{lstlisting}[style=rcode]
x_vals <- 0:2
y_vals <- 0:2

# Joint PMF matrix: rows = X (0,1,2), cols = Y (0,1,2)
P <- matrix(c(
  0.40, 0.10, 0.00,
  0.05, 0.25, 0.10,
  0.00, 0.05, 0.05
), nrow = 3, byrow = TRUE)

pX <- rowSums(P)
pY <- colSums(P)

indep_check_00 <- pX[1] * pY[1]  # Compare with P[1,1]

EX <- sum(x_vals * pX)
EY <- sum(y_vals * pY)

EXY <- sum(outer(x_vals, y_vals, "*") * P)
covXY <- EXY - EX * EY

EX2 <- sum((x_vals^2) * pX)
EY2 <- sum((y_vals^2) * pY)
varX <- EX2 - EX^2
varY <- EY2 - EY^2
rho <- covXY / sqrt(varX * varY)

list(pX = pX, pY = pY, independent = all.equal(P, outer(pX, pY)),
     indep_cell_00 = c(product = indep_check_00, p00 = P[1,1]),
     EX = EX, EY = EY, EXY = EXY, covXY = covXY, rho = rho)
\end{lstlisting}
\end{solutionbox}

\vspace{6pt}

\begin{problembox}{Problem 5}
Average waiting time: $\mu=10$ min, $\sigma^2=4$ min$^2$.
\end{problembox}

\begin{solutionbox}{Solution to Problem~5}
\[
\P(X\ge20)\le\frac{10}{20}=0.50,\qquad
\P(|X-10|\ge10)\le\frac{4}{100}=0.04.
\]
\textbf{R code:}
\begin{lstlisting}[style=rcode]
mu <- 10
var <- 4
a <- 20

# Markov bound: P(X >= a) <= E[X]/a
markov_bound <- mu / a

# Chebyshev bound for P(X >= 20):
# P(X >= 20) <= P(|X - mu| >= 10) <= var / 10^2
cheb_bound <- var / (20 - mu)^2

c(markov_bound = markov_bound, chebyshev_bound = cheb_bound)
\end{lstlisting}
\end{solutionbox}

\vspace{6pt}

\begin{problembox}{Problem 6}
A component is acceptable if length is $50\pm1$ mm, with $\mu=50$, $\sigma=0.2$.
\end{problembox}

\begin{solutionbox}{Solution to Problem~6}
\[
\frac{\delta}{\sigma}=\frac{1}{0.2}=5<6,\qquad
\sigma^*\le\frac{1}{6}\approx0.167.
\]
\textbf{R code:}
\begin{lstlisting}[style=rcode]
mu <- 50
sigma <- 0.2
L <- 49
U <- 51

delta <- min(U - mu, mu - L)
k_sigma <- delta / sigma
meets_6sigma <- k_sigma >= 6
sigma_required <- delta / 6

c(delta = delta, k_sigma = k_sigma, meets_6sigma = meets_6sigma,
  sigma_required = sigma_required)
\end{lstlisting}
\end{solutionbox}

\vspace{6pt}

\begin{problembox}{Problem 7}
For fair coin flips, find $n$ so that
$\P(|\bar X_n-0.5|\le0.01)\ge0.95$ via Chebyshev.
\end{problembox}

\begin{solutionbox}{Solution to Problem~7}
\[
\P(|\bar X_n-0.5|\ge0.01)\le\frac{2500}{n}\le0.05
\;\Rightarrow\;
n\ge50{,}000.
\]
\textbf{R code:}
\begin{lstlisting}[style=rcode]
# Chebyshev:
# P(|Xbar - 0.5| >= 0.01) <= Var(Xbar)/0.01^2 = (0.25/n)/0.0001 = 2500/n
n_min <- ceiling(2500 / 0.05)
n_min
\end{lstlisting}
\end{solutionbox}

\vspace{6pt}

\begin{problembox}{Problem 8}
Each of 100 boxes has $\mu=20$ kg and $\sigma=2$ kg.
Find $\P(\text{total}>2050)$.
\end{problembox}

\begin{solutionbox}{Solution to Problem~8}
\[
\P(S_{100}>2050)=\P\!\left(Z>\frac{2050-2000}{20}\right)=\P(Z>2.5)\approx0.0062.
\]
\textbf{R code:}
\begin{lstlisting}[style=rcode]
n <- 100
mu <- 20
sd <- 2
x <- 2050

mu_S <- n * mu
sd_S <- sqrt(n) * sd
p_over <- 1 - pnorm(x, mean = mu_S, sd = sd_S)

c(mu_S = mu_S, sd_S = sd_S, p_over = p_over)
\end{lstlisting}
\end{solutionbox}

\vspace{6pt}

\begin{problembox}{Problem 9}
Insurer has $n=10{,}000$ policies with payout:
$0$ (prob.\ $0.995$) or $\$50{,}000$ (prob.\ $0.005$).
\end{problembox}

\begin{solutionbox}{Solution to Problem~9}
\[
\begin{aligned}
\P(S_n>3{,}000{,}000)
&=\P\!\left(Z>\frac{3{,}000{,}000-2{,}500{,}000}{352{,}670}\right)\\
&=\P(Z>1.418)\approx0.078.
\end{aligned}
\]
\textbf{R code:}
\begin{lstlisting}[style=rcode]
n <- 10000
p_claim <- 0.005
claim_amount <- 50000
premium <- 300

# One-policy moments
mu <- p_claim * claim_amount
EX2 <- p_claim * claim_amount^2
varX <- EX2 - mu^2
sdX <- sqrt(varX)

# CLT approximation for total claims S_n
mu_S <- n * mu
sd_S <- sqrt(n) * sdX
R <- n * premium
z <- (R - mu_S) / sd_S
p_loss_clt <- 1 - pnorm(z)

# Exact (binomial) check:
# S_n = claim_amount * B, B ~ Binomial(n, p_claim)
# Net loss when S_n > R  <=>  B > R / claim_amount = 60
p_loss_exact <- 1 - pbinom(60, size = n, prob = p_claim)

c(mu = mu, sdX = sdX, mu_S = mu_S, sd_S = sd_S,
  p_loss_clt = p_loss_clt, p_loss_exact = p_loss_exact)
\end{lstlisting}
\end{solutionbox}

\bigskip
\section*{\color{kseblue}Bonus Problem: CLT (Easy Level)}

\begin{problembox}{Problem 10}
Average receipt is $\mu=\$12$, $\sigma=\$3$, with $n=100$ customers.
\end{problembox}

\begin{solutionbox}{Solution to Problem~10}
\[
\begin{aligned}
\P(S_{100}>1250)
&=\P\!\left(Z>\frac{1250-1200}{30}\right)\\
&=\P(Z>1.667)\approx0.0475.
\end{aligned}
\]
\textbf{R code:}
\begin{lstlisting}[style=rcode]
n <- 100
mu <- 12
sd <- 3

mu_S <- n * mu
sd_S <- sqrt(n) * sd

# (b) P(S_100 > 1250)
p_over_1250 <- 1 - pnorm(1250, mean = mu_S, sd = sd_S)

# (c) 95th percentile v where P(S_100 > v) = 0.05
v95 <- qnorm(0.95, mean = mu_S, sd = sd_S)

c(mu_S = mu_S, sd_S = sd_S, p_over_1250 = p_over_1250, v95 = v95)
\end{lstlisting}
\end{solutionbox}

\end{document}
